Given an $S_n$-character $\phi$, let  
\begin{equation}\label{e_k_interpretation1}
    m_{k,\phi}
    := \langle \phi,\chi^{(n-k,1^k)} \rangle 
    \quad (0 \le k \le n-1)
    , \quad
    e_{k,\phi} 
    :=\langle \phi,\chi^{(1^k)\oplus(n-k)} \rangle
    \quad (0 \le k \le n).
\end{equation}
When $\phi$ is understood from the context, we use the abbreviated notations $m_k:=m_{k,\phi}$ and $e_k:=e_{k,\phi}$.

By Pieri's rule~\cite[Theorem 7.5.17]{EC2} 
combined with the (inverse) Frobenius characteristic map,
\begin{align*}
    \chi^{(1^k)\oplus (n-k)} 
    = \chi^{(1^k)} \chi^{(n-k)}
    &= \ch^{-1}(s_{(1^k)}s_{(n-k)})
    = \ch^{-1}(s_{(n-k+1,1^{k-1})} + s_{(n-k,1^k)})\\
    &= \chi^{(n-k+1,1^{k-1})} + \chi^{(n-k,1^k)} 
    \qquad (0<k<n).
\end{align*}
Equivalently,
\[
    \chi^{(n-k,1^k)}=
    \sum_{i=0}^{k} (-1)^{k-i}\chi^{(1^i)\oplus (n-i)} \qquad (0\le k \le n-1).
\]
Thus the sequences $\{m_k\}_{k=0}^{n-1}$ and $\{e_k\}_{k=0}^{n}$ 
determine each other via the relations
\begin{equation}\label{eq:alternating}
    e_k = m_k + m_{k-1}
    \quad \text{and} \quad
    m_k = \sum_{i=0}^k (-1)^{k-i} e_i
    \qquad (0 \le k \le n),
\end{equation}
where $m_k := 0$ for $k = -1$ and $k = n$.
Note that, in particular,
\[
    \sum_{i=0}^{n} (-1)^{n-i} e_i = 0.
\]


\subsection{Cyclic descent extension and hook-multiplicities}\label{sec:hook_criterion}
	
Let $\A\subseteq S_n$ be Schur-positive with $\Q(\A) = \ch(\phi)$.
Denote
\[
    m_\lambda 
    := \langle \phi, \chi^\lambda \rangle 
    = \langle \Q(\A), s_\lambda \rangle
\qquad (\forall \lambda \vdash n).
\]
For $\lambda = (n-k,1^k)$ we use the abbreviation 
\[
m_k := m_{(n-k,1^k)} 
\qquad (0 \le k < n).
\]

The {\em hook-multiplicity generating function} is defined as 
\[
M_{\A}(x):=\sum_{k=0}^{n-1} m_k x^k.
\]

The function $M_{\A}(x)$, where $\A$ 
is a conjugacy class, 
was studied and applied to the enumeration of unimodal permutations with a given cycle type by Thibon~\cite{Thibon}. 

\begin{observation}\label{lem:1}
    For every $0\le k <n$
    \[
        m_k = \card{\{a\in \A:\ \Des(a)=[k]\}}.
    \]
\end{observation}

\begin{proof}
    For every $0 \le k < n$ there exists 
    a unique standard Young tableau $T$ of size~$n$ with
    $\Des(T) = [k]$ (where $[0] := \varnothing$). 
    The shape of $T$ is $(n-k,1^k)$.
    Comparing the coefficients of ${\bf x}^{[k]}$ on both sides of Equation~\eqref{eq:l2} completes the proof.
\end{proof}

We now restate and prove Lemma~\ref{lem:21}.

\begin{lemma}\label{lem:2}
    A Schur-positive set $\A \subseteq S_n$ carries a cyclic descent extension if and only if 
    the hook-multiplicity generating function $M_\A(x)$ 
    is divisible by $1+x$
    and the quotient $M_\A(x)/(1+x)$ has non-negative coefficients; 
    equivalently, if and only if
    there exist non-negative integers $d_k$ $(0 \le k \le n-2)$ such that
    \[
        m_k = d_k + d_{k-1} \qquad (0 \le k \le n-1),
    \]
    where $d_k := 0$ for $k = -1$ and $k = n-1$.
\end{lemma}
	
\begin{proof}
    If $\A$ carries a cyclic descent extension then, by Observation~\ref{lem:1} and the equivariance of $\cDes$, for every $0 \le k \le n-1$:
    \begin{align*}
        m_k
        &= \card{\{a\in \A:\ \Des(a)=[k]\}} \\
        &= \card{\{a\in \A:\ \cDes(a)=[k]\}}+\card{\{a\in \A:\ \cDes(a)=[k]\sqcup\{n\}\}} \\
        &= \card{\{a\in \A:\ \cDes(a)=[k]\}}+\card{\{a\in \A:\ \cDes(a)=[k+1]\}}.
    \end{align*}
    The numbers
    \[
        d_k := \card{\{a\in \A:\ \cDes(a)=[k+1]\}}
        \qquad (-1 \le k \le n-1)
    \]
    satisfy the required conditions,
    which imply the corresponding properties of $M_\A(x)$.

    For the opposite direction, assume that
    there exist non-negative integers $d_k$ (with $d_{-1} = d_{n-1} = 0$)
    such that $m_k=d_{k-1}+d_k$ for all $0\le k\le n-1$. 
    By Pieri's rule~\cite[Theorem 7.15.7]{EC2}, 
    \[
        s_{(1^k) \oplus (n-k)}=s_{(1^k)} s_{(n-k)}=s_{(n-k+1,1^{k-1})}+s_{(n-k,1^k)}
        \qquad (1 \le k \le n-1).
    \]
    Hence
    \[
        \sum_{k=1}^{n-1} d_{k-1} s_{(1^k)\oplus (n-k)} 
        = \sum_{k=1}^{n-1} d_{k-1} (s_{(n-k+1,1^{k-1})}+s_{(n-k,1^k)})
        =\sum_{k=0}^{n-1} m_k s_{(n-k,1^k)}. 
    \]
    Since
    \[
        \Q(\A) 
        = \sum_{\lambda \vdash n} \langle \Q(\A), s_\lambda \rangle s_\lambda 
        = \sum_{\substack{\lambda\vdash n \\ \lambda \ \text{non-hook}}} m_\lambda s_\lambda +
        \sum_{k=0}^{n-1} m_k s_{(n-k,1^k)}
    \]
    we obtain
    \begin{equation}\label{eq:q}
        \Q(\A) 
        = \sum_{\substack{\lambda\vdash n \\ \lambda \ \text{non-hook}}} m_\lambda s_\lambda +
        \sum_{k=1}^{n-1} d_{k-1} s_{(1^k)\oplus (n-k)}.
    \end{equation}
    By Corollary~\ref{cor:Schur-gf}, this is equivalent to
    \[
        \sum_{a\in \A} {\bf x}^{\Des(a)} = \sum_{\substack{\lambda\vdash n \\ \lambda\ \text{non-hook}}} m_\lambda 
        \sum_{T\in \SYT(\lambda)} {\bf x}^{\Des(T)} + 	\sum_{k=1}^{n-1} d_{k-1}
        \sum_{T \in \SYT((1^k)\oplus (n-k))} {\bf x}^{\Des(T)}.
    \]
    By Theorem~\ref{conj1},  
    the set $\SYT(\lambda)$ carries a cyclic descent extension if and only if $\lambda \vdash n$ is not a hook; 
    and each of the sets $\SYT((1^k)\oplus (n-k))$ $(1 \le k \le n-1)$ carries a cyclic descent extension. 
    Hence $\A$ also carries a cyclic descent extension, completing the proof. 
\end{proof}

\begin{corollary}
    If a Schur-positive set $\A$ carries a cyclic descent extension then 
    \[
        \sum_{a\in \A} {\bf x}^{\cDes(a)} = \sum_{\substack{\lambda\vdash n \\ \lambda\ \text{non-hook}}} m_\lambda 
        \sum_{T\in \SYT(\lambda)} {\bf x}^{\cDes(T)} + 	\sum_{k=1}^{n-1} d_{k-1}
        \sum_{T\in \SYT((1^k)\oplus (n-k))}{\bf x}^{\cDes(T)},
    \]
    where $m_\lambda$ and $d_k$ are the non-negative integers defined above.
\end{corollary}

\begin{proof}
    By Corollary~\ref{cor:cDes-dist-Schur}
	together with Equation~\eqref{eq:q}, 
	the generating function for the corresponding cyclic descent set is uniquely determined	and satisfies 
    the claimed equality.
\end{proof}
	
	
\subsection{Divisibility of the hook-multiplicity generating  function}
\label{sec:divisibility}	
	
Recall the notation 
\[
    m_{k,\phi} :=
    \langle \phi,\chi^{(n-k,1^k)} \rangle 
    \qquad (0 \le k < n)
\]
for an $S_n$-character $\phi$.

\begin{lemma}\label{lem:div}
    For every $S_n$-character $\phi$, the hook-multiplicity generating function 
    \[
    M_\phi(x) 
    := \sum\limits_{k=0}^{n-1} m_{k,\phi} x^k
    \]
    is divisible by $1+x$ if and only if 
    the value of $\phi$ on an $n$-cycle is zero, i.e.,  $\phi_{(n)}=0$. 
\end{lemma}

\begin{proof}
    By~\cite[Lemma 4.10.3]{Sagan},  
    for every partition $\lambda\vdash n$ 
    \[
        \chi^\lambda_{(n)} =
        \begin{cases} 
            (-1)^k,&\text{if } \lambda = (n-k,1^k) \text{ for some } 0 \le k < n;\\
		  0,&\text{otherwise.}
		\end{cases}
    \]
    Thus   
    \[
        \phi_{(n)}
        = \sum\limits_{\lambda\vdash n} \langle \phi,\chi^\lambda \rangle \chi^\lambda_{(n)}
        = \sum\limits_{k=0}^{n-1} \langle \phi,\chi^{(n-k,1^k)} \rangle \chi^{(n-k,1^k)}_{(n)}
        = \sum\limits_{k=0}^{n-1} m_{k,\phi} \cdot (-1)^k
        = M_\phi(-1),
    \]
    which equals zero if and only if $1+x$ divides $ M_\phi(x)$, 
    completing the proof.
\end{proof}

Letting $\phi = \psi^\lambda$, the higher Lie character indexed by a partition $\lambda$, reduces Proposition~\ref{prop:divisibility} to the following character evaluation. 

Recall the M\"obius function $\mu(n)$, the sum of all primitive (complex) $n$-th roots of $1$. If $n$ has a prime square divisor then $\mu(n)=0$; otherwise, $n$ is a product of $k$ distinct primes and $\mu(n)=(-1)^k$. 
The following lemma is equivalent to a combinatorial identity due to Garsia, as shown in  Proposition~\ref{prop:equiv_Garsia} below. We give here an independent direct algebraic proof.

\begin{lemma}\label{lem:char_eval_cycles}
    For $\lambda \vdash n$
    \[
        \psi^\lambda_{(n)} =
        \begin{cases} 
            \mu(r), &\text{if } \lambda = (r^s);\\
	       0, &\text{otherwise,}\\
        \end{cases}
    \]
    where $\mu$ is the M\"obius function.
\end{lemma}

\begin{proof}
    Let $c$ be an $n$-cycle in $S_n$, and let  $Z_\lambda = Z_{S_n}(g)$ be the centralizer in $S_n$ of a specific element $g\in \C_\lambda$. 
    An explicit formula for the induced character~\cite[(5.1)]{Isaacs} is
    \[
        \psi^\lambda(c)=\omega^\lambda\uparrow_{Z_\lambda}^{S_n}(c)=\frac{1}{|Z_\lambda|}\sum_{\substack{x\in S_n\\x^{-1}cx\in Z_\lambda}}\omega^\lambda(x^{-1}cx).
    \]
    An $n$-cycle commutes only with its own powers. Thus, if $\lambda$ is {\em not} of the form $(r^s)$ for some $r$ and $s$, 
    then there is no $n$-cycle in $Z_\lambda$;
    equivalently, $x^{-1}cx \not\in Z_\lambda$ for every~$x\in S_n$. 
    It follows that, for such partitions $\lambda \vdash n$, $\psi^\lambda_{(n)}=0$.

    Assume now that $\lambda= (r^s)$, and let let $g = g_1 g_2 \cdots g_s \in \C_\lambda$ be a fixed  product of~$s$ disjoint $r$-cycles. 
    The order of the centralizer 
    $Z_\lambda = Z_{S_n}(g)$ is $s!r^s$. 
    If $u \in Z_\lambda$ is an $n$-cycle ($n = rs$) then 
    $g = u^k$ for some integer $k$ with $\gcd(k,n) = s$;
    equivalently, $u^s = g^j$ for some $0 < j < r$ with $\gcd(j,r) = 1$.
    Conversely, if $u \in S_n$ is an $n$-cycle satisfying 
    $u^s = g^j$ for some $0 < j < r$ with $\gcd(j,r) = 1$, then $g$ is a power of $u$ and therefore $u \in Z_\lambda$.
    Thus the number of $n$-cycles in $Z_\lambda$, namely the number of elements of $Z_\lambda \cap \C_{(n)}$, is $\varphi(r) (s-1)! r^{s-1}$, where $\varphi$ is Euler's totient function. 

    Viewing $Z_\lambda$ as the group of $s \times s$ monomial (``generalized permutation'') matrices whose nonzero entries are complex $r$-th roots of unity, an element $u \in Z_\lambda \cap \C_{(n)}$ corresponds to a matrix whose underlying permutation is a full $s$-cycle and the product of its nonzero entries is a primitive $r$-th root of unity.
    This product is equal to $\omega^\lambda(u)$, so it is a primitive $r$-th root of unity.
    For $u, v \in Z_\lambda \cap \C_{(n)}$, write $u \sim v$ if $v = u^i$ for some integer $i$ (necessarily coprime to $n$). This clearly defines an equivalence relation on $Z_\lambda \cap \C_{(n)}$. On each equivalence class, all primitive $r$-th roots of unity appear with the same frequency as values of $\omega^\lambda$. This property thus holds for all of $Z_\lambda \cap \C_{(n)}$, where this frequency is $(s-1)!r^{s-1}$. Denoting by $\xi$ any specific primitive $r$-th root of unity, the sum of all values of $\omega^\lambda$ on $Z_\lambda \cap \C_{(n)}$ is therefore
    \[
        \sum_{u\in Z_\lambda\cap\C_{(n)}} \omega^\lambda(u)
        = (s-1)! r^{s-1} \sum_{j:(j,r)=1} \xi^j
        = (s-1)! r^{s-1} \mu(r).
    \]
    Given any $c,u \in \C_{(n)}$, there are exactly $n=rs$ permutations $x \in S_n$ which satisfy $u=x^{-1}cx$.
    Thus
    \begin{align*}
        \psi^\lambda(c)
        &= \omega^\lambda\uparrow_{Z_\lambda}^{S_n}(c)
        = \frac{1}{|Z_\lambda|}\sum_{\substack{x \in S_n \\ x^{-1}cx \in Z_\lambda}} \omega^\lambda(x^{-1}cx)
        = \frac{n}{|Z_\lambda|}\sum_{u \in Z_\lambda\cap\C_{(n)}} \omega^\lambda(u) \\
        &= \frac{n}{s!r^s} (s-1)! r^{s-1} \mu(r)
        = \mu(r), 
    \end{align*}
    as claimed.
\end{proof}

\begin{proof}[Proof of Proposition~\ref{prop:divisibility}.]
	By Lemma~\ref{lem:div}, 
	$1+x$ divides the hook-multiplicity generating function of the higher Lie character $\psi^\lambda$ 
    if and only if $\psi^\lambda_{(n)}=0$.
    Lemma~\ref{lem:char_eval_cycles} completes the proof.
\end{proof}

\begin{corollary}\label{cor:only if}
    Let $\lambda \vdash n$.
    \begin{enumerate}
    \item
        If $\lambda = (r^s)$ for some square-free integer $r$ and positive integer $s$, then $1+x$ does not divide the hook-multiplicity generating function $M_\lambda(x)$, and the descent set map on the conjugacy class $\C_\lambda$ does not have a cyclic extension.
    \item 
        If $\lambda$ is not equal to $(r^s)$ for any square-free $r$, then $1+x$ divides $M_\lambda(x)$. 
        In this case, the descent set map on 
        $\C_\lambda$ has a cyclic extension if and only if 
        the quotient $M_\lambda(x)/(1+x)$ has non-negative coefficients.
    \end{enumerate}
\end{corollary}

\begin{proof}
    By the Gessel--Reutenauer theorem (Theorem~\ref{thm:GR22}), for every $\lambda\vdash n$ the conjugacy class $\C_\lambda$ 
    is Schur-positive, with $\Q(\C_\lambda) = \ch(\psi^\lambda)$.
    Combining this with Lemma~\ref{lem:2}  and  Proposition~\ref{prop:divisibility} completes the proof of both parts.
\end{proof}

In the following sections we will prove the non-negativity of the coefficients of the quotient $M_\lambda(x)/(1+x)$ for partitions (cycle types) which are not equal to $(r^s)$ for a square-free $r$: 
cycle types with more than one cycle length will be considered in Section~\ref{sec:distinct}, 
single cycles will be considered in Section~\ref{sec:cycles},
and cycle types $\lambda = (r^s)$ with non square-free $r$ and $s>1$ will be considered in Section~\ref{sec:Witt} (this is the most difficult case).


\section{Non-negativity: the case of more than one cycle length}
\label{sec:distinct}

Consider, first, the case of a conjugacy class with more than one cycle length. This is the easiest case to handle.
